\section*{الفصل \ref{ch:b1:radiometric-methods} --- طرائق التأريخ الإشعاعي}
% القيم: calc/ch25.py

\begin{solution}{exo:b1:radiometric-methods:1}
(أ)~$D^*/P = 2^n - 1 = 3$ يعطي $n = 2$: نصفا عمر. (ب)~الروبيديوم-87:
$e^{\lambda t} - 1 = e^{\num{1.3972e-11} \times \num{3.10e8}} - 1 = 0.004341$؛
اليورانيوم-235: $e^{\num{9.8485e-10} \times \num{3.10e8}} - 1 = 0.3570$. ونسبة اليورانيوم أكبر بنحو 80 مرة: فخطأ قياس مطلق معيّن جزء أصغر بكثير منها. (ج)~$2^{-n} < 0.01$ حين $n > \log_2 100 = 6.64$: بعد نحو سبعة أنصاف أعمار.
\end{solution}

\begin{solution}{exo:b1:radiometric-methods:2}
(أ)~كربون مشع على الفحم. (ب)~يورانيوم–رصاص على الزركون، أو أرغون–أرغون (بوتاسيوم–أرغون) على السانيدين، وهو \omterm{def:b1:rock-forming-minerals:mineral}{معدن} غني بالبوتاسيوم.
(ج)~يورانيوم–رصاص على الزركون، الذي يصمد أمام تاريخ طويل؛ والروبيديوم–السترونتيوم على صخور كاملة ممكن؛ أما البوتاسيوم–الأرغون فيُخاطر بأعمار خفضها أرغون فُقد في تسخين لاحق. (د)~الكربون المشع على كربونات الصدفة، قرب حدّ الطريقة (نحو \num{70000} سنة). (ه)~لا شيء مباشرة: عمرها محصور بين صخور بركانية مؤرّخة فوقها وتحتها، أو يُعرف بترابط حفرياتها مع قطاع مؤرّخ (\cref{ch:b1:stratigraphic-scale}).
\end{solution}

\begin{solution}{exo:b1:radiometric-methods:3}
(أ)~$t = -8033 \ln 0.25 = 8033 \times 1.3863 = \num{11136}$ سنة BP.
(ب)~متوسط العمر هو $5700/\ln 2 = \num{8223}$ سنة: $t = 8223 \times 1.3863 =
\num{11400}$ سنة. (ج)~كلاهما يفترض أن نسبة الكربون-14 في الكربون الحيّ كانت آنذاك كما في المعيار؛ وهذا لم يكن كذلك. ويُقرأ العمر التقويمي من منحنى معايرة مثل IntCal20، مدًى من السنين cal BP.
\end{solution}

\begin{solution}{exo:b1:radiometric-methods:4}
(أ)~تبدأ كل ساعة حين يبرد معدنها دون درجة إغلاقه:
الزركون (فوق \qty{750}{\celsius}) يُغلق أولًا، ثم الهورنبلند (نحو \qty{530}{\celsius})، ثم البيوتايت (نحو \qty{280}{\celsius}) أخيرًا، فتتناقص الأعمار بهذا الترتيب مع برودة الغرانيت. (ب)~$(530 - 280)/(410 - 398) =
250/12 = \qty{21}{\celsius}$ لكل مليون سنة. (ج)~أغلق الزركون فوق
\qty{750}{\celsius}، فبرد الغرانيت بمقدار $750 - 530 =
\qty{220}{\celsius}$ على الأقل في \qty{10}{\Myr}: أي \qty{22}{\celsius} على الأقل لكل مليون سنة، وهذا حدّ أدنى.
\end{solution}

\begin{solution}{exo:b1:radiometric-methods:5}
(أ)~$\bar x = 7.5$، $\bar y = 0.809525$، $S_{xx} = 81.0$، $S_{xy} = 1.13965$؛
$m = 1.13965/81.0 = 0.0140698$، $b = 0.809525 - 7.5 \times 0.0140698 = 0.70400$.
(ب)~$t = \ln(1.0140698)/\num{1.3972e-11} = \qty{1000.0}{\Ma}$؛ \omterm{def:b1:radiometric-methods:isochron}{النسبة الابتدائية}
$0.7040$. (ج)~البواقي هي $+1.6$، و$-2.5$، و$+0.7$ و$+0.2$ (بوحدات
\num{e-4}): $s = \sqrt{\num{9.34e-8}/2} = 0.00022$، و$\sigma_m = 0.00022/\sqrt{81}
= \num{2.40e-5}$، و$\sigma_t = \num{2.40e-5}/(\num{1.3972e-11} \times 1.0141) =
\qty{1.7}{\Myr}$: \qty[uncertainty-mode=separate]{1000.0(17)}{\Ma}. (د)~$0.7040$
قريب من قيمة \omterm{def:b1:layered-earth:layers}{الوشاح} $0.7025$ وبعيد جدًّا عن $0.710$ القارية: فالمصدر في معظمه في \omterm{def:b1:layered-earth:layers}{الوشاح}.
\end{solution}

\begin{solution}{exo:b1:radiometric-methods:6}
(أ)~$0.100/39.098 \times \num{1.17e-4} = \qty{2.992e-7}{mol}$ من البوتاسيوم-40 لكل
غرام. (ب)~${}^{40}\mathrm{Ar}^*/{}^{40}\mathrm{K} = \num{1.169e-9}/\num{2.992e-7}
= 0.0039065$؛ $\lambda/\lambda_e = 5.543/0.581 = 9.5404$؛ $t = \ln(1 + 9.5404
\times 0.0039065)/\num{5.543e-10} = \ln(1.037269)/\num{5.543e-10} =
\qty{66.01}{\Ma}$. (ج)~الأرغون الإشعاعي سيكون $0.995 \times 1.169$ نانومول،
فيعطي \qty{65.69}{\Ma}: فالعمر المقاس سيكون أقدم من الواقع ب\qty{0.32}{\Myr}، أي عشرة أضعاف \num{32000} سنة التي يتزامن فيها الاصطدام والانقراض عند نهاية الطباشيري الآن. والأرغون الزائد يقلق أساسًا الصخور الفتية والصخور الكاملة؛ \omterm{def:b1:rock-forming-minerals:mineral}{ومعدن} مثل السانيدين لا ينبغي أن يحبس الأرغون حين يتبلّر، وطيف أعمار أرغون–أرغون كان سيكشفه.
\end{solution}

\begin{solution}{exo:b1:radiometric-methods:7}
(أ)~نحو \num{2000} cal BP: يمرّ المنحنى هناك ب\num{2070} سنة BP. (ب)~من نحو \num{2495} إلى \num{2699} cal BP، أي مائتا سنة: فالمنحنى شبه مسطّح بين \num{2500} و\num{2700} cal BP، لذلك يكون لكل عمر تقويمي في هذا المدى عمر كربون مشع شبه واحد، ولن يضيّق قياس أدقّ المدى. (ج)~عند \num{2000} cal BP يكون \omterm{def:b1:radiometric-methods:radiocarbon}{عمر الكربون المشع} \emph{أقدم} ب\num{70} سنة من العمر التقويمي؛ وعند \num{8000} cal BP أصغر ب\num{822} سنة (\qty{10}{\percent})؛ وعند \num{12000} cal BP أصغر ب\num{1694} سنة (\qty{14}{\percent}). والفرق ليس متناسبًا مع العمر: بل يغيّر إشارته. ولا يمكن أن يأتي من نصف العمر وحده، الذي كان سيعطي \qty{2.4}{\percent} ثابتة.
\end{solution}

\begin{solution}{exo:b1:radiometric-methods:8}
(أ)~الرماد تحت الطبقة (\qty{94.50}{\Ma}) يعطي حدًّا أقصى، وتدفّق الحمم فوقها (\qty{93.80}{\Ma}) حدًّا أدنى: فالحدّ محصور بينهما. والدايك حدّ أدنى أيضًا، لكنه أضعف، وأحدث من الحمم. والزركون الفُتاتي حدّ أقصى للحجر الرملي وحده، وأضعف من الرماد الذي يقع تحت الحجر الرملي. (ب)~$f = (12 - 4)/(20 - 4) = 0.5$؛ $t = 94.50 - 0.5 \times 0.70 =
\qty{94.15}{\Ma}$؛ $\sigma_t = \sqrt{0.25 \times 0.10^2 + 0.25 \times 0.12^2} =
\qty{0.08}{\Myr}$: \qty[uncertainty-mode=separate]{94.15(8)}{\Ma}. (ج)~جاءت حبيباته من صخر عمره \qty{95.3}{\Ma} على الأقل، تآكل في مكان ما من منطقة المصدر؛ وهذا يبيّن أن الحجر الرملي أحدث من \qty{95.3}{\Ma}، وهو ما يثبته الرماد السفلي بإحكام أكبر.
\end{solution}

\begin{solution}{exo:b1:radiometric-methods:9}
(أ)~الحبّة~A: $\ln(1.08401)/\num{1.55125e-10} = \qty{520.0}{\Ma}$ و
$\ln(1.6688)/\num{9.8485e-10} = \qty{520.0}{\Ma}$. الحبّة~B: \qty{368.4}{\Ma} و
\qty{389.9}{\Ma}. (ب)~الحبّة~A متوافقة. (ج)~$0.05881/0.08401 = 0.7000$ و
$0.4682/0.6688 = 0.7001$. والعمر هو $\ln(1 + R)/\lambda$، وليس متناسبًا مع
$R$: فلنسبة اليورانيوم-238 الصغيرة يكون $\ln(1+R) \approx R$، فينخفض العمر بنحو \qty{30}{\percent} (إلى \qty{368}{\Ma})؛ أما لنسبة اليورانيوم-235 الأكبر فينخفض اللوغاريتم بأقلّ، ولا يهبط العمر إلا إلى \qty{390}{\Ma}. وكلاهما أحدث من اللازم، وهما مختلفان: فالحبّة غير متوافقة.
\end{solution}

\begin{solution}{exo:b1:radiometric-methods:10}
(أ)~مع $u = \lambda t$، يكون $D^*/P = e^u - 1 = u + u^2/2 + \dots$، ف$(D^*/P)/\lambda
= t\,(1 + u/2 + \dots)$: يُبالغ التقريب في العمر بخطأ نسبي قريب من $u/2$. (ب)~الروبيديوم-87 عند \qty{310}{\Ma}: $u = 0.00433$، الخطأ
\qty{0.22}{\percent} (\qty{0.7}{\Myr})؛ البوتاسيوم-40 عند \qty{296}{\Ma}: $u =
0.164$، الخطأ \qty{8.7}{\percent} (\qty{26}{\Myr}؛ و$u/2$ يعطي \qty{8.2}{\percent})؛
اليورانيوم-238 عند \qty{4.4}{\Ga}: $u = 0.683$، الخطأ \qty{43}{\percent}، فالتقريب وحتى تقدير خطئه ($u/2 = \qty{34}{\percent}$) يفشلان. (ج)~$R = 0.004341$، $\sigma_R = \num{4.34e-6}$؛ وبموجب
\cref{prop:b1:radiometric-methods:precision}، $\sigma_t =
\num{4.34e-6}/(\num{1.3972e-11} \times 1.0043) = \qty{0.31}{\Myr}$.
\end{solution}

\begin{solution}{exo:b1:radiometric-methods:11}
(أ)~لتكن ${}^{40}\mathrm{K}$ البوتاسيوم-40 اليوم. عند الزمن $t_2$ قبل الحاضر كانت هناك ${}^{40}\mathrm{K}\,e^{\lambda t_2}$ ذرّة، والأرغون المصنوع بين $t_1$ و$t_2$ هو، بموجب \cref{prop:b1:radiometric-methods:kar} مطبّقة على ذلك الفاصل، $(\lambda_e/\lambda)\,{}^{40}\mathrm{K}\,e^{\lambda
t_2}\bigl(e^{\lambda (t_1 - t_2)} - 1\bigr) = (\lambda_e/\lambda)\,
{}^{40}\mathrm{K}\,\bigl(e^{\lambda t_1} - e^{\lambda t_2}\bigr)$. وبقي منه كسر
$1 - f$. ومنذ $t_2$، أنتج \omterm{def:b1:rock-forming-minerals:mineral}{المعدن} المغلق
$(\lambda_e/\lambda)\,{}^{40}\mathrm{K}\,(e^{\lambda t_2} - 1)$ زيادة أخرى. وبجمع الأجزاء والضرب في $\lambda/(\lambda_e\,{}^{40}\mathrm{K})$ نحصل على التعبير.
(ب)~العمر الظاهري هو $\ln(1 + \text{التعبير})/\lambda$: \qty{400.0}{\Myr} من أجل
$f = 0$، و\qty{342.3}{\Myr} من أجل $0.3$، و\qty{282.7}{\Myr} من أجل $0.6$ و
\qty{200.0}{\Myr} من أجل $1$. (ج)~كل عمر بين \qty{200}{\Myr} و
\qty{400}{\Myr} يمكن أن ينتجه فقدان جزئي ما، ولا يستطيع رقم واحد أن يقول إن كان يسجّل حدثًا حقيقيًّا في ذلك الزمن أم خليطًا؛ فنحتاج إلى طيف أعمار، أو إلى ساعات أخرى في الصخر نفسه.
\end{solution}

\begin{solution}{exo:b1:radiometric-methods:12}
(أ)~العمر يتناسب مع متوسط العمر المستعمل: $10306 \times 8223/8033 =
\num{10550}$ سنة. (ب)~$10550 - 10306 = 244$ سنة، أي نحو سُبع ال\num{1694}. (ج)~إن كان النشاط الابتدائي $k A_0$، فالعمر الحقيقي هو $t = \tau
\ln(k A_0/A) = \num{10550} + \tau \ln k$؛ ومع $t = \num{12000}$، يكون $\ln k =
1450/8223 = 0.176$ و$k = 1.19$: أي نحو \qty{19}{\percent} من الكربون-14 أكثر مما في المعيار.
\end{solution}

\begin{solution}{pb:b1:radiometric-methods:1}
\begin{enumerate}
  \item تصعد النقاط بانتظام مع $x$؛ وانحرافها عن خطّ لا يتجاوز وحدتين في الرقم العشري الرابع، مقابل مدى قدره $0.04$ في $y$: فهي على خطّ مستقيم مع تشتّت صغير.
  \item $\bar x = 4.90333$، $\bar y = 0.729288$، $S_{xx} = 62.1253$، $S_{xy} =
        0.269335$؛ $m = 0.269335/62.1253 = 0.0043354$.
  \item $b = 0.729288 - 0.0043354 \times 4.90333 = 0.70803$: نسبة
        $^{87}\mathrm{Sr}/^{86}\mathrm{Sr}$ للغرانيت حين تبلّر، مشتركة بين كل أجزائه.
  \item $t = \ln(1.0043354)/\num{1.3972e-11} = \qty{309.6}{\Ma}$.
  \item البواقي هي $+9.1$، و$-17.2$، و$+6.9$، و$-6.4$، و$+15.1$ و$-7.5$ (بوحدات
        \num{e-5})؛ $s = 0.000137$، $\sigma_m = 0.000137/\sqrt{62.1253} =
        \num{1.74e-5}$، $\sigma_t = \num{1.74e-5}/(\num{1.3972e-11} \times 1.0043)
        = \qty{1.2}{\Myr}$: \qty[uncertainty-mode=separate]{309.6(12)}{\Ma}، و
        $\sigma_b = 0.00010$.
  \item $0.7080$ أقرب إلى $0.710$ القارية منه إلى $0.7025$ \omterm{def:b1:layered-earth:layers}{للوشاح}.
        وللخليط، $(0.70803 - 0.7025)/(0.710 - 0.7025) = 0.74$: نحو ثلاثة أرباع
        السترونتيوم جاءت من \omterm{def:b1:layered-earth:layers}{قشرة} منصهرة.
  \item $0.0740/39.098 \times \num{1.17e-4} = \qty{2.214e-7}{mol}$ لكل غرام،
        مضروبة في $\num{6.022e23}$: \num{1.33e17} ذرّة.
  \item $^{40}\mathrm{Ar}^*/^{40}\mathrm{K} = \num{4.139e-9}/\num{2.2144e-7} =
        0.018691$؛ $1 + 9.5404 \times 0.018691 = 1.17832$؛ $t =
        \ln(1.17832)/\num{5.543e-10} = \qty{296.0}{\Ma}$.
  \item $^{40}\mathrm{K} = \qty{3.292e-8}{mol.g^{-1}}$، والنسبة $0.019239$،
        و$\ln(1.18355)/\num{5.543e-10} = \qty{304.0}{\Ma}$.
  \item $\ln(1.018691)/\num{5.543e-10} = \qty{33.4}{\Ma}$، أحدث بتسع مرّات، لأن
        واحدًا من كل عشرة اضمحلالات تقريبًا يعطي أرغونًا. والكالسيوم-40 هو النظير الشائع
        للكالسيوم، الوفير في الصخور: فالقليل الذي يصنعه الاضمحلال لا يُميَّز عن الكالسيوم
        الذي احتوته \omterm{def:b1:rock-forming-minerals:mineral}{المعادن} منذ البداية.
  \item الأرغون الزائد يزيح الأعمار بقدر نحو مليون سنة، وهذا يهمّ الصخور الفتية؛ وهنا سيكون أقلّ من \qty{1}{\percent} من
        العمر. \omterm{def:b1:rock-forming-minerals:groups}{والميكا} \omterm{def:b1:rock-forming-minerals:groups}{والأمفيبول} تحبس قليلًا جدًّا من الأرغون حين تتكوّن.
  \item إما أن الغرانيت بَرد ببطء فأغلقت \omterm{def:b1:rock-forming-minerals:groups}{الميكا} \omterm{def:b1:rock-forming-minerals:groups}{والأمفيبول} بعد التبلّر (\omterm{def:b1:radiometric-methods:closure}{أعمار تبريد})، أو أن تسخينًا لاحقًا طرد جزءًا من أرغونها. ويدعم هضبة مستوية في طيف أعمار أرغون–أرغون عمرَ التبريد؛ وطيف بلا هضبة يحذّر من أن \omterm{def:b1:rock-forming-minerals:mineral}{المعادن} فقدت الأرغون. وترتيب الأعمار، الهورنبلند قبل البيوتايت كما تتنبأ به درجات إغلاقها، يدلّ أيضًا على التبريد.
  \item $\ln(1.049361)/\num{1.55125e-10} = \qty{310.6}{\Ma}$ و
        $\ln(1.35784)/\num{9.8485e-10} = \qty{310.6}{\Ma}$: متوافقان.
  \item الزركون \qty{310.6}{\Ma}، وإغلاقه فوق \qty{750}{\celsius}: التبلّر؛ \omterm{def:b1:radiometric-methods:isochron}{والإيزوكرون} للصخر الكامل،
        \qty[uncertainty-mode=separate]{309.6(12)}{\Ma}، يؤرّخه أيضًا، ويتفق ضمن عدم يقينه؛ والهورنبلند \qty{304.0}{\Ma}: التبريد
        عبر نحو \qty{530}{\celsius}؛ والبيوتايت \qty{296.0}{\Ma}: التبريد عبر نحو
        \qty{280}{\celsius}. ينحدر المنحنى من فوق \qty{750}{\celsius} عند \qty{310.6}{\Ma} إلى \qty{280}{\celsius} عند \qty{296.0}{\Ma}.
  \item $(530 - 280)/(304.02 - 296.03) = 250/7.99 = \qty{31}{\celsius}$ لكل
        مليون سنة.
  \item $(750 - 530)/(310.60 - 304.02) = 220/6.58 = \qty{33}{\celsius}$ لكل
        مليون سنة على الأقل، لأن الزركون أُغلق فوق
        \qty{750}{\celsius}: فكان التبريد على الأقل بالسرعة نفسها في البداية.
  \item فرق الدرجات، $250 \pm \sqrt{40^2 + 40^2} = 250 \pm
        57$، يعطي \qtyrange{24}{38}{\celsius} لكل مليون سنة.
  \item $(550 - 350)/7.99 = \qty{25}{\celsius}$ لكل مليون سنة، داخل المدى
        في السؤال~17: فعدم اليقين يأتي من \omterm{def:b1:radiometric-methods:closure}{درجات الإغلاق}، لا من الأعمار.
  \item $(280 - 10)/30 = \qty{9.0}{km}$.
  \item $\ln(1.039108)/\num{1.55125e-10} = \qty{247.30}{\Ma}$، والمعطى نفسه من نسبة
        اليورانيوم-235: متوافقان.
  \item من الغرانيت، المتآكل عند \omterm{def:b1:cross-sections:unconformity}{عدم التوافق}: فعمرهما هو عمر تبلّره. وهما
        يدلّان فقط على أن الحجر الرملي أحدث من \qty{310}{\Ma}: حدّ أقصى، أضعف بكثير من
        طبقة الرماد.
  \item $(58 - 18)/(247.30 - 245.30) = \qty{20.0}{m}$ لكل مليون سنة؛ وقاعدة الحجر الرملي
        أقدم من الرماد السفلي ب$18/20.0 = \qty{0.90}{\Myr}$:
        \qty{248.20}{\Ma}.
  \item من \qty{309.6}{\Ma} (\omterm{def:b1:radiometric-methods:isochron}{الإيزوكرون}) أو \qty{310.6}{\Ma} (الزركون) إلى
        \qty{248.2}{\Ma}: نحو \qty{62}{\Myr}.
  \item \qty{9.0}{km} في $296.0 - 248.2 = \qty{47.8}{\Myr}$: \qty{0.19}{km} لكل
        مليون سنة، أي \qty{0.19}{mm} في السنة. وهذا يفترض أن البيوتايت المؤرَّخ جاء من موضع
        أسفل \omterm{def:b1:cross-sections:unconformity}{عدم التوافق} مباشرة، فأُزيل كل ما فوقه في الغرانيت والشيست بحلول وصول الرمل، وأن
        \omterm{def:b1:heat-budget:geotherm}{المنحنى الحراري} بقي ثابتًا.
  \item برد الغرانيت بنحو \qty{31}{\celsius} لكل مليون سنة بين إغلاقي الهورنبلند والبيوتايت.
\end{enumerate}
\end{solution}
